\section{Conclusion}
\label{sec:conclusion}

We evaluated GPT-4o zero-shot Gherkin generation from user stories on $N{=}100$ pairs sampled from the Rathnayake et al.~\cite{rathnayake2026llm} industrial corpus across four IntelligenceBank domains, using pre-registered automated metrics and one-tailed statistical tests.

We investigated whether the median cosine similarity between generated and expert-written scenarios exceeds 0.85 when measured with OpenAI \texttt{text-embedding-3-large} embeddings (RQ1).
Results show that the semantic similarity threshold was not achieved on the full sample (median$=0.7742$, $p \approx 1.000$, Cliff's $\delta=-0.70$): zero-shot GPT-4o outputs remain semantically divergent from expert references under this metric, even when they read plausibly on the surface.

We investigated whether the proportion of generated scenarios that parse as executable Gherkin exceeds 80\% (RQ2).
Results show that the executable syntax threshold was achieved (rate$=100.00\%$, $p=2.04\times10^{-10}$, Cohen's $h=0.93$), confirming that GPT-4o consistently produces structurally valid BDD scenarios under our parser-based definition.

Overall, GPT-4o \emph{failed} RQ1 but \emph{passed} RQ2: generated scenarios are syntactically usable yet semantically misaligned with expert references at the pre-defined similarity gate.
Practitioners should treat model output as a draft requiring human review against domain-specific acceptance criteria.

This is the first study to evaluate GPT-4o zero-shot Gherkin generation using automated metrics with statistical significance testing and fixed, pre-registered thresholds on a replicated industrial sample.
Our open replication package at \url{https://github.com/khanhkhanh55355/RT-SWT-003-nhom3} documents the full pipeline from sampling through API logging to hypothesis testing.

Future work should pursue three concrete directions.
First, evaluate \textbf{few-shot} and \textbf{chain-of-thought} prompting on the same $N{=}100$ dataset and thresholds to determine whether in-context guidance improves semantic alignment without eroding syntax compliance.
Second, extend evaluation beyond single happy-path scenarios to \textbf{multi-step flows}, negative paths, and boundary conditions that better reflect real backlog complexity.
Third, investigate \textbf{domain-specific fine-tuning} for Gherkin generation to assess whether semantic similarity can reach the 0.85 threshold while preserving the high executable syntax rate observed here.
